\chapter*{الملخص}
\addcontentsline{toc}{chapter}{الملخص}
\setlength{\parindent}{0pt}
\setlength{\parskip}{0.8em}

تعتمد المدن الحديثة اعتماداً كبيراً على شبكات واسعة من كاميرات المراقبة لأغراض الأمن ورصد حركة المرور وضمان السلامة العامة. غير أن معظم الأنظمة القائمة تعجز عن تتبع الأشخاص والمركبات بصورة مستمرة حين ينتقلون بين كاميرات مختلفة، ولا سيما عندما لا تتداخل مجالات رؤية هذه الكاميرات. ويؤدي ذلك إلى فقدان الهوية بصورة متكررة، وتجزؤ المسارات، وضعف القدرة على تحليل أنماط الحركة في المناطق الواسعة.

يعالج هذا المشروع مشكلة تتبع الكائنات عبر كاميرات متعددة مع الحفاظ على هوياتها بمرور الوقت. ويتمثل الهدف الرئيسي في تصميم وتنفيذ نظام قادر على اكتشاف الكائنات محل الاهتمام، وتتبعها داخل نطاق كل كاميرا، ثم ربط مظاهرها المرئية عبر الكاميرات المختلفة بطريقة موثوقة. ويهدف النظام إلى دعم سيناريوهات المراقبة على مستوى المدينة، حيث يتم توزيع عدد كبير من الكاميرات في مواقع مختلفة ذات زوايا رؤية وظروف إضاءة متباينة.

يعتمد النهج المقترح على خط معالجة يجمع بين اكتشاف الكائنات، والتتبع داخل الكاميرا الواحدة، وإعادة التعرف عبر الكاميرات. في البداية، يتم اكتشاف الكائنات محل الاهتمام في كل إطار من تدفقات الفيديو. بعد ذلك، يقوم مكوّن التتبع بربط هذه الاكتشافات زمنياً لبناء مسارات متصلة داخل كل كاميرا. ثم يقارن مكوّن إعادة التعرف السمات المرئية المستخرجة من كاميرات مختلفة لتحديد ما إذا كانت الكائنات التي ظهرت في مشاهد منفصلة تعود إلى الكيان الحقيقي نفسه. كما تم تصميم خط معالجة بيانات للتعامل مع مدخلات الفيديو، وإدارة النتائج الوسيطة، وتخزين معلومات التتبع النهائية لاستخدامها لاحقاً في التحليل.

ينتج عن المشروع نموذج أولي عملي يوضح إمكانية تتبع الكائنات عبر كاميرات متعددة مع الحفاظ على الهوية في مقاطع فيديو مسجلة. وقد تم تقييم النظام باستخدام سيناريوهات اختبار مختارة للتحقق من قدرته على الحفاظ على الهوية بشكل متسق عند انتقال الكائنات بين الكاميرات، والتعامل مع التحديات العملية مثل الحجب وتغيرات المظهر. وقد استُخدمت في التطوير أدوات معيارية في الرؤية الحاسوبية والتعلم الآلي، وتمثل النتائج خطوة أولى نحو منصة قابلة للتوسع لتطبيقات المراقبة الذكية على مستوى المدينة.